%!TEX program = xelatex
%!TEX encoding = ISO-8859-1
\documentclass[12pt]{article}

\usepackage{amssymb}
\usepackage{amsmath,amsfonts,amsthm,amstext}
\usepackage[brazil]{babel}
%\usepackage[latin1]{inputenc}
\usepackage{graphicx}
\graphicspath{{/home/jfreitas/Dropbox/imagens-latex/}{/Users/jfreitas/Dropbox/imagens-latex/}}
\usepackage{enumitem}
\usepackage{multicol}
\usepackage[all]{xy}

\setlength{\topmargin}{-1.0in}
\setlength{\oddsidemargin}{0in}
\setlength{\textheight}{10.1in}
\setlength{\textwidth}{6.5in}
\setlength{\baselineskip}{12mm}

\newcounter{exercicios}
\setcounter{exercicios}{0}
\newcommand{\questao}{
\addtocounter{exercicios}{1}
\noindent{\bf Exerc{\'\i}cio \arabic{exercicios}: }}

\newcommand{\equi}{\Leftrightarrow}
\newcommand{\bic}{\leftrightarrow}
\newcommand{\cond}{\rightarrow}
\newcommand{\impl}{\Rightarrow}
\newcommand{\nao}{\sim}
\newcommand{\sub}{\subseteq}
\newcommand{\e}{\ \wedge\ }
\newcommand{\ou}{\ \vee\ }
\newcommand{\vaz}{\emptyset}

\newcommand{\real}{\mathbb{R}}
\newcommand{\vesp}{\vspace{0.2cm}}
\newcommand{\z}{\mathbb{Z}}
\newcommand{\n}{\mathbb{N}}
\newcommand{\q}{\mathbb{Q}}
\newtheorem{defin}{Defini{\c c}{\~a}o}

\newcommand{\compcent}[1]{\vcenter{\hbox{$#1\circ$}}}
\newcommand{\comp}{\mathbin{\mathchoice
{\compcent\scriptstyle}{\compcent\scriptstyle}
{\compcent\scriptscriptstyle}{\compcent\scriptscriptstyle}}}

\begin{document}
\pagestyle{empty}

\begin{figure}[h]
        \begin{minipage}[c]{1.7cm}
        \includegraphics[width=1.7cm]{unb.pdf}
        \end{minipage}%
        \hspace{0pt}
        \begin{minipage}[c]{4in}
          {Universidade de Bras{\'\i}lia} \\
          {Departamento de Matem{\'a}tica}
\end{minipage}
\end{figure}
\vspace{-1cm}\hrule

\begin{center}
{\Large\bf {\'A}lgebra 1 - Turma A} \\ \vspace{9pt} {\large\bf
  $3^{\underline{a}}$ Lista de Exerc{\'\i}cios -- 2$^{o}$/2011}\\ \vspace{9pt} Prof. Jos{\'e} Ant{\^o}nio O. Freitas
\end{center}
\hrule

\begin{center}
\Large{\bf An{\'e}is}
\end{center}

\vspace{.6cm}

\questao Ache todos os ideais do anel $\z_4$, do anel $\z_{12}$ e do anel
$\z_{16}$.

\vesp

\questao Mostre que a interse{\c c}{\~a}o de quaisquer dois ideais de um anel
comutativo {\'e} sempre um conjunto n{\~a}o vazio. Mostre que esse conjunto n{\~a}o
vazio {\'e} um ideal tamb{\'e}m.

\vesp

\questao Ache os ideais primos dos an{\'e}is $\z_2$, $\z_4$, $\z_5$,
$\z_8$ e $\z_{10}$. O ideal $\{\overline{0}\}$ do anel $\z_4$ {\'e} primo?

\vesp

\questao Seja $I$ um ideal de $\z$. Mostre que se todos os elementos n{\~a}o
nulos de $\z/I$ s{\~a}o invers{\'\i}veis, ent{\~a}o $I$ {\'e} um ideal primo.

\vesp

\questao As seguintes afirma{\c c}{\~o}es s{\~a}o verdadeiras ou falsas?
\begin{multicols}{2}
\begin{enumerate}[label={\alph*})]
\item $2\z \sub 4\z$;
\item$3\z \sub 7\z$;
\item$8\z \sub 2\z$;
\item$2\z \cap 8\z = 8\z$;
\item$3\z \cap 7\z = 21\z$;
\item$n\z \cap m\z = mn\z$;
\item$2\z \cup 3\z = 6\z$;
\item$2\z \cup 8\z = 2\z$;
\item$n\z \cup m\z = mdc(m,n)\z$.
\end{enumerate}
\end{multicols}

\vesp

\questao Seja $p$ um n{\'u}mero primo. Seja $A$ um ideal de $\z$, tal que $p\z
\sub A$. Mostre que $A = \z$ ou $A = p\z$.

\vesp

\questao Sejam $m$, $n$ inteiros positivos. Mostre que a interse{\c c}{\~a}o $m\z
\cap n\z$ {\'e} um ideal de $\z$. {\'E} poss{\'\i}vel achar o gerador desse ideal?

\vesp

\questao {\'E} verdadeiro ou falso que a uni{\~a}o de dois ideais de $\z$ {\'e} um ideal
de $\z$?

\vesp

\questao Sejam $A$ e $B$ dois ideais de $\z$. Definimos a {\bf soma de
  ideais}
\[
A + B = \{ a + b \mid a \in A, b \in B\}
\]
quando $a$ varia sobre todos os elementos de $A$ e $b$ varia sobre todos os
elementos de $B$. Mostre que a soma de dois ideais de $\z$ {\'e} um ideal de $\z$.

\vesp

\begin{defin}
Seja $(A, + , \cdot)$ um anel. Dizemos que um subconjunto n{\~a}o vazio $C\subseteq A$ {\'e} um \textbf{subanel} de $A$ quando $(C, + ,\cdot)$ {\'e} um anel.
\end{defin}

\vesp

\questao Seja $(A, + , \cdot)$ um anel. Mostre que subconjunto n{\~a}o vazio $C\subseteq A$ {\'e} um \textbf{subanel} de $A$ se, e somente se, $x - y \in C$ e $x\cdot y \in C$ para todos $x$, $y \in C$.


\newpage

\questao Determinar quais dos seguintes subconjuntos de $\q$ s{\~a}o suban{\'e}is:
\begin{multicols}{2}
\begin{enumerate}[label=({\alph*})]
\item $\z$
\item $B = \{x \in \q \mid x \notin \z\}$
\item $C = \left\{\dfrac{a}{b} \in \q \mid a \in \z,\ b \in \z,\ 2 |b \right\}$
\item $D = \left\{\dfrac{a}{2^n} \in \q \mid a \in \z \mbox{ e } n \in \z \right\}$
\end{enumerate}
\end{multicols}
\vesp

\questao Verifique se $L = \{ a + b\sqrt{2} \mid a, b \in \q\}$ {\'e} um subanel
do anel $\mathbb{R}$.

\vesp

\questao Considere os an{\'e}is $\z$ e $\z\times \z$. Verique se s{\~a}o homomorfismos:
\begin{enumerate}[label=({\alph*})]
\item $f : \z\times\z \to \z\times\z$ dado por $f(x,y) = (0,y)$
\item $f : \z\times\z \to \z$ dado por $f(x,y) = y$
\item $f : \z\to \z\times\z$ dado por $f(x) = (2x,0)$
\item $f : \z\times\z \to \z\times\z$ dado por $f(x,y) = (-y,-x)$
\item $f : \z \to \z\times\z$ dado por $f(x) = (0,x)$
\end{enumerate}

\vesp

\questao Seja $f: A \to B$ um homomorfismo de an{\'e}is. Mostre que:
\begin{enumerate}[label=({\alph*})]
\item Se $C$  {\'e} um subanel de $A$, ent{\~a}o $f(C)$ {\'e} um subanel de $B$.
\item Se $D$ {\'e} um subanel de $B$, ent{\~a}o $f^{-1}(D)$ {\'e} um subanel de $A$.
\item Se $I$ {\'e} um ideal de $A$, ent{\~a}o $f(I)$ {\'e} um ideal de $B$.
\item Se $J$ {\'e} um ideal de $B$, ent{\~a}o $f^{-1}(J)$ {\'e} um ideal de $A$.
\end{enumerate}

\vesp

\questao D{\^e} um exemplo de an{\'e}is $A$ e $B$ e um homomorfismo $f : A \to B$ tal que $f(1_A) \ne 1_B$.

\vesp

%\questao Sejam $(A, +, \cdot)$ e $(B, \oplus, \otimes)$ dois an{\'e}is e $f: A
%\to B$ um homomorfismo. Definimos o {\it n{\'u}cleo} ou o {\it
%  kernel} de $f$ como o conjunto
%  \[
%  Ker(f) = \{a \in A \mid f(a) = 0_B\}.
%  \]
%Mostre que o kernel {\'e} um ideal de $A$.

\vesp

\questao Mostre que o kernel de $f$ pode ser considerado como a imagem
inversa do subconjunto $\{0_B\}$ pela fun{\c c}{\~a}o $f$.

{\it Sugest{\~a}o: Seja $f: X \to Y$ uma fun{\c c}{\~a}o}. A imagem inversa {\it
  de um subconjunto $U \sub Y$ {\'e} o subconjunto $X$ definido por $f^{-1}(U) = \{x \in X \mid f(x) \in U\}$. }

\vesp

%\questao Mostre que um homomorfismo do anel $A$ no anel $B$ {\'e} um
%monomorfismo se, e somente se, $Ker(f) = \{0_A\}$.

\vesp

\questao Determine o kernel dos homomorfismos do Exerc{\'i}cio 13.

\vesp

\questao Sejam os an{\'e}is $A = \{ a + b\sqrt{-2} \mid a,\ b \in \q\}$ e $B = M_2(\q)$.
\begin{enumerate}[label=({\alph*})]
\item Mostre que $f : A \to B$ dada por
\[
f(a + b\sqrt{-2}) =
\begin{pmatrix}
a & -2b\\
b & a
\end{pmatrix}
\]
{\'e} um homomorfismo.
\item $f$ {\'e} um isomorfimo?
\end{enumerate}

\questao {\'E} verdadeiro ou falso: $\z$ e $\z_{m}$ para $m > 1$ s{\~a}o an{\'e}is
isomorfos.

\vesp

\questao Considere os seguintes an{\'e}is: $(\real, +, \cdot)$ e $(\real, \oplus, \odot)$, sendo $a \oplus b = a + b + 1$ e $a \odot b = a + b + ab$. Mostre que $f : \real \to \real$ dado por $f(x) = x + 1$, para todo $x \in \real$, {\'e} um isomorfimos de $(\real, \oplus, \odot)$ em $(\real, +, \cdot)$. Defina o isomorfismo inverso.

\vesp

\questao Mostre que nenhuma aplica{\c c}{\~a}o $f : A \to B$, onde $A = \{ x + y\sqrt{2} \mid x,\ y \in \q\}$ e $B = \{ x + y\sqrt{3} \mid x,\ y \in \q\}$ {\'e} isomorfismo.

{\it Sugest{\~a}o: Observe que se $f$ fosse um isomorfimo de $A$ em $B$, ent{\~a}o $f(\sqrt{2}) = a + b\sqrt{3}$. Calcule a seguir $f(2) = 2$ a partir de $f(\sqrt{2}) = a + b\sqrt{3}$}

\vesp

\questao Seja $I$ um ideal de um anel $A$. Mostre que se $A$ possui unidade, ent{\~a}o $A/I$ tamb{\'e}m possui unidade.

\vesp

\questao Suponha que $A$ {\'e} um anel com unidade e $I$ um ideal de $A$. Mostre que $a + I \in A/I$ {\'e} invers{\'\i}vel se, e somente se, existe $r \in A$ de modo que $ar - 1 \in I$.

\vesp

\questao D{\^e} um exemplo de um anel de integridade $A$ e de um ideal $I$ em $A$ tal que $A/I$ n{\~a}o {\'e} de integridade.

\vesp

\begin{defin}
Um \textbf{corpo} {\'e} um anel associativo, comutativo e unit{\'a}rio, no qual todos os elementos n{\~a}o nulos s{\~a}o invers{\'i}veis, isto {\'e}, possuem inverso na multiplica{\c c}{\~a}o.
\end{defin}

\vesp

\questao Seja $A$ um corpo. Mostre que o inverso de cada elemento {\'e} {\'u}nico. Em outras palavras, um elemento n{\~a}o nulo n{\~a}o pode ter mais de um inverso.

\vesp

\questao Mostre que todo corpo tem pelo menos dois elementos.

\vesp

\questao Mostre que todo corpo {\'e} anel de integridade.

\vesp

\questao Se $A$ {\'e} um corpo, mostre que seus {\'u}nicos ideais s{\~a}o os triviais.

\vspace{0.6cm}

\hrule
\begin{center}
\Large{\bf Grupos}
\end{center}
\hrule

\vspace{0.6cm}

\questao Mostre que num grupo $G$, o elemento neutro e inverso de cada
elemento {\'e} {\'u}nico.

\vesp

\questao Verifique se o conjunto $\q_{>0}$ dos n{\'u}meros racionais estritamente positivos com a
 opera{\c c}{\~a}o dada {\'e} ou n{\~a}o um grupo. Justifique sua
resposta.
\begin{multicols}{2}
\begin{enumerate}[label=({\alph*})]
\item $(\q_{>0},\cdot)$
\item $(\q_{>0}, +)$
\end{enumerate}
\end{multicols}

\vesp

\questao Seja $z  = a + bi \in \mathbb{C}$, onde $a$, $b \in \real$. Definimos $|z| = \sqrt{a^2 + b^2}$. Prove que $G=\{z \in \mathbb{C} \mid |z| = 1\}$ {\'e} um grupo
abeliano com a opera{\c c}{\~a}o de multiplica{\c c}{\~a}o de n{\'u}meros complexos.

\vesp

\questao Mostre que o conjunto $\q[\sqrt{2}]^*=\{ a + b\sqrt{2} \in
\mathbb{R}^* \mid  a, b \in \q \}$ {\'e} um grupo multiplicativo abeliano.

\vesp

\questao Seja $(G,*)$ um grupo com elemento neutro $e$. Para $x\in
G$, considere a nota{\c c}{\~a}o $x^n=x*x*\cdots *x$ ($n$ vezes).
\begin{enumerate}[label=({\alph*})]
\item Seja $G$ um grupo tendo $e$ como elemento neutro. Prove que se
$x^2=e$, para todo $x\in G$, ent{\~a}o $G$ {\'e} um grupo abeliano.
\item Mostre que se $x\in G$ {\'e} tal que $x^2=x$, ent{\~a}o $x$ {\'e} o elemento neutro.
\end{enumerate}

\vesp

\questao Seja $(G, *)$ um grupo e $g\in G$ um elemento fixado e $g^{-1}$
seu inverso. Mostre que a aplica{\c c}{\~a}o $i_g: G\to G$ definida por
$i_g(x)=g^{-1}*x*g$, para todo $x \in G$, {\'e} um isomorfismo.

\newpage

\questao Quais dos seguintes subconjuntos $G$ de $\z_{13}$ s{\~a}o grupos
com a opera{\c c}{\~a}o de multiplica{\c c}{\~a}o?
\begin{multicols}{2}
\begin{enumerate}[label=({\alph*})]
\item $G=\{\overline{1},\overline{12}\}$;

\item $G=\{\overline{1},\overline{5},\overline{8},\overline{12}\}$;

\item $G=\{\overline{1},\overline{2},\overline{3},\overline{4}, \overline{5},\overline{6},\overline{7},
 \overline{8},\overline{9},\overline{10},\overline{11},\overline{12}\}$
\item $G=\{\overline{1}, \overline{3},\overline{5},\overline{7},\overline{9},\overline{11}\}$.
\end{enumerate}
\end{multicols}

\vesp

\questao Sejam $H$ e $K$ subgrupos de um grupo $G$ (com nota{\c c}{\~a}o
multiplicativa).
\begin{enumerate}[label=({\alph*})]
\item Mostre que $H\cap K$ tamb{\'e}m {\'e} subgrupo de $G$.
\item Seja $g\in G$ um elemento fixado. Mostre que o conjunto
$g^{-1}Hg=\{ g^{-1}xg \mid x\in H \} $ {\'e} um subgrupo de $G$.
\item Prove que $H\cup K$ {\'e} subgrupo de $G$ se, e somente se,
$H\subseteq K$ ou $K\subseteq H$.
\item Demonstre que $HK=\{hk \mid h\in H, k\in K\}$ {\'e} subgrupo
de $G$ se, e somente se, $HK=KH$.

[\emph{Nota: $HK=KH$ \textbf{n{\~a}o} quer dizer que $hk=kh$,
para todo $h\in H, k\in K$; significa que $hk=k_1h_1 \in KH$ e $kh=h_2k_2 \in
HK$, para todo $h\in H, k\in K$.}]
\end{enumerate}

\vesp

\questao Sabendo que $G = \{e, a, b, c, d, f\}$ {\'e} um grupo multiplicativo isomorfo ao grupo $(\z_6, +)$, pede-se:
\begin{enumerate}[label=({\alph*})]
\item construir uma tabela de multiplica{\c c}{\~a}o para $G$;
\item calcular $a^2$, $b^{-2}$ e $c^{-3}$;
\item obter $x \in G$ tal que $dxf = a^{-1}$.
\end{enumerate}

\vesp

\questao Sejam $f:G\to H$ e $g:H\to J$ isomorfismos de grupos. Mostre
que
\begin{enumerate}[label=({\alph*})]
\item $g\circ f$ {\'e} tamb{\'e}m isomorfismo de grupos;
\item $f^{-1}$ {\'e} isomorfismo de grupos.
\end{enumerate}

\vesp

\questao Sejam $G$ e $J$ grupos multiplicativos, $f : G \to J$ um homomorfismo de grupos e $H$ um subgrupo de $J$. Mostre que $f^{-1}(H) = \{ x \in G \mid f(x) \in H\}$ {\'e} um subgrupo de $G$.

\vesp

\questao Prove que um grupo $G$ {\'e} abeliano se, e somente se, $f : G \to G$ definada por $f(x) = x^{-1}$ {\'e} um homomorfismo.

\vesp

\questao Seja $f: G\to H$ um
homomorfismo sobrejetivo de grupos e $K$ um subgrupo de $H$.
\begin{enumerate}[label=({\alph*})]
\item Mostre que $f^{-1}(K)$ {\'e} um subgrupo de $G$, onde $f^{-1}(K)$ {\'e} a imagem inversa de $K$.
\item Mostre que o n{\'u}cleo Ker$(f)$ de $f$ {\'e} um subgrupo de
$G$  e que Ker$(f)\subset f^{-1}(K)$.
\item Seja $N$ um subgrupo de $G$. Mostre que $\mbox{Ker}(f)\cdot
N = f^{-1}(f(N))$, onde $f^{-1}(f(N))$ {\'e} a imagem inversa de $f(N)$.

[Lembre que $\mbox{Ker}(f)\cdot
N=\{xy ~|~ x\in \mbox{Ker}(f), y\in N\}$ e mostre que $\mbox{Ker}(f)\cdot
N \subseteq f^{-1}(f(N))$. Para ver que $f^{-1}(f(N))\subseteq \mbox{Ker}(f)\cdot
N$, tome $x\in f^{-1}(f(N))$. Ent{\~a}o $f(x)\in f(N)$, ou seja, $f(x)=f(y)$, para algum $y\in N$. Assim,
$f(xy^{-1})=e$ e portanto $xy^{-1}=z \in \mbox{Ker}(f)$. Ent{\~a}o $x=zy \in \mbox{Ker}(f)\cdot
N$.]
\end{enumerate}

\vesp

\questao Seja $f: G \to J$ um homomorfismo de grupos e $g\in G$ tal que $o(g)=n$.
\begin{enumerate}[label=({\alph*})]
\item Mostre que $f(g)$ tem ordem positiva e que a ordem de $f(g)$ divide $n$.
\item Mostre que se $f$ {\'e} isomorfismo, ent{\~a}o $o(f(g))=n$.
\end{enumerate}

\vesp

\questao Seja $G$ um grupo finito com elemento
neutro $e$ e suponha que $H$ e $K$ s{\~a}o subgrupos de $G$ tais que
$|H|=p$ e $|K|=q$, onde $p$ e $q$ s{\~a}o primos distintos. Mostre que
$H\cap K=\{e\}$.

\end{document}